\begin{proposicion}[Reconstruction of the inner class]\label{meso:reconstruccion}
There are constants $C_{\rm int},\eta_{\rm int},\gamma_{\rm int}>0$, depending only on the profile, with the
following property. Suppose the observables of this section, namely the proportions in each interval, the
fibers, and the sixteen rectangle terms, are within their tolerances, and suppose
$D_{\max}(k)<\alpha_{\rm geo} k$ with $\alpha_{\rm geo}=\varepsilon/(4\|W_0\|_\infty)$. If the mesh
$a_{\rm mes}$, the tolerance $\varrho_{\rm mes}$, and $\varepsilon$ are sufficiently small and $k$ is
large, the Bethe point of the board exists, satisfies the diffusion, vacancy, and gap bounds with strict
inequalities for these constants, and has mass less than $\varepsilon$ on each diagonal. On the available
squares $p_e/q_e=1+O(\varrho_{\rm mes}+\varepsilon+1/k)$, with a constant depending on the profile.
Moreover, for these values of $\varepsilon$ and large $k$, every child of a board whose Bethe point
satisfies these bounds with diagonal mass at most $\varepsilon$ satisfies the outer bounds, and each of
its diagonals has mass at most $\varepsilon_{\rm ext}=2\varepsilon$.
\end{proposicion}
\begin{proof}
By the fiber bounds, each row and column of $q$ sums to $1+O(\varrho_{\rm mes})$. For $V$, approximate
$\nu_1(r,\cdot)$ or $\nu_1(\cdot,c)$ by a constant on each interval. The continuous normalization of the
profile, its Lipschitz regularity, and the sum of the interval tolerances give degrees $1+O(a_{\rm
mes}+\varrho_{\rm mes})$ uniformly in the coordinate of the tracked line. Fix
${0<\theta<\min(c_\nu,1)/2}$. The matrix $V-\theta\mathbf1\mathbf1^T/k$ is nonnegative and has row and
column sums ${\le1-\theta+O(a_{\rm mes}+\varrho_{\rm mes})}$. By Schur's inequality its operator norm is
at most the same quantity. The rank-one approximation and the bound on $q-V$ give
\[
 \sigma_2(q)\le1-\theta+O(a_{\rm mes}+\varrho_{\rm mes})+\|q-V\|_{\rm op}.
\]
If $D_R,D_C$ are the diagonal matrices of degrees of $q$, then $D_R^{-1/2}qD_C^{-1/2}$ differs from $q$ by
$O(\varrho_{\rm mes})$ in operator norm. With the preceding tolerances sufficiently small, its second
singular value is $\le1-\theta/2$. Write $L_q$ for the weighted row--column Laplacian,
$D=\operatorname{diag}(D_R,D_C)$ and $d_{\min}=\min_v D_{vv}$. The normalized adjacency has
off-diagonal blocks $D_R^{-1/2}qD_C^{-1/2}$ and its transpose, so its eigenvalues are the signed
singular values. Its top eigenvalue is one, with eigenvector $D^{1/2}\mathbf1$. The normalized
Laplacian $D^{-1/2}L_qD^{-1/2}$ thus has gap
$\lambda=1-\sigma_2(D_R^{-1/2}qD_C^{-1/2})\ge\theta/2$. For every real vector $f$, its Rayleigh
inequality after subtracting the degree-weighted mean gives
\[
 f^TL_q f\ \ge\ \lambda\min_{a\in\mathbb R}\sum_v D_{vv}(f_v-a)^2
 \ \ge\ \lambda d_{\min}\min_{a\in\mathbb R}\sum_v(f_v-a)^2 .
\]
For $f\perp\mathbf1$ the last minimum is $\|f\|_2^2$, so the gap of $L_q$ is at least
$d_{\min}\lambda$.
Here all degrees are $1+O(\varrho_{\rm mes})$.

We detail the passage from these observables to the Bethe point. On the support of $q$, let $B$ be the
row--column incidence matrix of Section~\ref{macro:reconstruccion} and let $P_q=\operatorname{diag}(q_e)$.
For rebalancing, solve $(BP_qB^T)z=\mathbf1-Bq$ on the orthogonal complement of the gauge
$(\mathbf1,-\mathbf1)$. The right-hand side is orthogonal to that vector. Changing the signs of the column
coordinates turns the matrix into the row--column Laplacian with weights $q$. The gap obtained above
and the upper and lower degree bounds give $\|z\|_2/\sqrt{k}=O(\varrho_{\rm mes})$. In the equation
for each vertex, the contribution from its neighbors is at most
$(\sum_j q_{ij}^2)^{1/2}\|z\|_2=O(\varrho_{\rm mes})$, since $q_{ij}\le\|W_0\|_\infty/k$, and
likewise for columns. Dividing by the degree gives $\|z\|_\infty=O(\varrho_{\rm mes})$. Hence
\[
 \bar q=q+P_qB^Tz,\qquad B\bar q=\mathbf1,\qquad
 \bar q_e/q_e=1+O(\varrho_{\rm mes}).
\]
Comparison of the weights shows that the diffusion and gap bounds still hold with strict inequalities.
Since $D_{\max}(k)<\alpha_{\rm geo} k$ with $\alpha_{\rm geo}=\varepsilon/(4\|W_0\|_\infty)$, each
diagonal in either family has
$q$-mass at most $\varepsilon/4$ and $\bar q$-mass at most $\varepsilon/2$ for $\varrho_{\rm mes}$
sufficiently small independently of $\varepsilon$.

To make the inversion in relative variables precise, set $P=\operatorname{diag}(\bar q)$ and
\[
 \mathcal T=\{x:BPx=0\},\qquad
 \Pi=I-B^T(BPB^T)^+BP ,
\]
where the superscript $+$ denotes the inverse on the orthogonal complement of the gauge, extended by zero.
This projection onto $\mathcal T$ annihilates $\operatorname{im}B^T$, the space of row and column
multipliers, and has uniformly bounded $\ell^\infty$ norm. Indeed, the preceding Laplacian estimate
applies with right-hand side $BPv$, whose norm is $\le\|v\|_\infty$ because $B\bar q=\mathbf1$. In
the feasible variables ${p=\bar q\odot(1+x)}$, $x\in\mathcal T$, the KKT equation is
$F(x)=-\Pi\nabla H_B(p)=0$, where $H_B$ is the Bethe entropy \eqref{bethe}, with the board as its
subscript rather than the incidence matrix.
Modulo row and column multipliers, its gradient has components $-\log p_e-3\log(1-p_e)+\log h_{D(e)}+\log
h_{A(e)}$. At $\bar q$, we have $\log q\in\operatorname{im}B^T$ because $W_0=b_0 r_Rr_C$, and $\log(\bar
q/q)=B^Tz+O(\varrho_{\rm mes}^2)$ in the uniform norm. The terms from the vacancies and from $1-\bar q_e$
are $O(\varepsilon+1/k)$. Consequently $\|F(0)\|_\infty\le C(\varrho_{\rm mes}^2+\varepsilon+1/k)$.

On the ball $\|x\|_\infty\le t<1/2$, the derivative of $\log p$ in relative variables is the diagonal
matrix with entries $1/(1+x_e)$, which differs from the identity by $O(t)$. The derivative of
$\log(1-p_e)$ has norm $O(1/k)$, and each derivative of $\log h_f$ has norm $\le\sum_{e\in f}\bar
q_e/h_f=O(\varepsilon)$. The projection is the identity on $\mathcal T$, so $\|DF(x)-I\|_{\infty,\mathcal
T}\le C(t+\varepsilon+1/k)$. Fix $t$ small and then take $\varrho_{\rm mes},\varepsilon,1/K_0$
sufficiently small that $x\mapsto x-F(x)$ is a contraction with constant $\le1/2$ on the ball and
$2\|F(0)\|_\infty\le t$. The fixed point satisfies
\[
 \|p/\bar q-1\|_\infty\le 2\|F(0)\|_\infty,\qquad
 p_e/q_e=1+O(\varrho_{\rm mes}+\varepsilon+1/k).
\]
All ratios are taken on available squares. The point is the unique stationary point on the convex domain
$0<p_e<1/4$, $h_f>0$, since there the negative Hessian of the Bethe entropy is the sum of a positive
diagonal matrix with entries $1/p_e-3/(1-p_e)$ and positive semidefinite forms from the lines. Since
$p_{\max}\le C/k\le e^{-4}$ for large $k$, the scalar calculation in Proposition~\ref{macro:maximo} also
shows that the point constructed is the unique global maximizer. It therefore lies on the branch selected
in the representation of \cite[Theorem~8.1 and the paragraph following it]{pred}. We obtain strict
diffusion, gap, and vacancy bounds, with diagonal mass $<\varepsilon$. The construction used only
observables of the board reached. For every child of a last regular parent, even if one of these
observables fails its bound, the inheritance result of \cite[Sections~6.5 and~8.2--8.3]{pred}, which
is item (I1) of Section~\ref{leyes}, gives, on the squares that remain available, a relative variation
$1+O(1/k)$ with respect to the parent, with constants from the outer class. The diagonal mass of the
child is at most $(1+O(1/k))\varepsilon\le\varepsilon_{\rm ext}=2\varepsilon$ after increasing $K_0$.
\end{proof}
